% !TEX root = ../main.tex
\section{Methods}

\hypertarget{case-study-the-umpqua-national-forest-forplan-model}{%
\subsection{Case study: Umpqua National Forest FORPLAN model}\label{case-study-the-umpqua-national-forest-forplan-model}}

The case study reproduces the Umpqua National Forest FORPLAN planning model \citep{usda1987umpqua}, which was used in the thesis, reconstructed in the open-source \texttt{ws3} wood-supply model \citep{paradis2026ws3}. The forest is stratified into four ecoclasses of decreasing growth potential: Douglas-fir-dominated \textit{CH-CW}, hemlock-dominated \textit{CD-CP}, fir-dominated \textit{CR-CF}, and hemlock-dominated \textit{CM-CE}. Seven silvicultural prescriptions are available, each with ecoclass-specific rotation-age ranges. The planning horizon is 150~years, divided into 15~periods of 10~years. The objective is maximization of net present value at a 4\% discount rate (varied in the experiments), with delivered-log prices escalating at 1\% per year for the first 50 years and constant thereafter. Volumes are expressed in thousands of cubic feet (MCF). The harvest-flow policy is the bounded-deviation sequential-flow constraint of equation~\eqref{eq:openloop_flow}. The scenarios include six policy levels, represented by different harvest-flow constraint sets (Section~\ref{experiment-design}) and a set of landbase conditions, i.e.. initial forest; the landbase of central interest (landbase~1) is an all-mature forest, the extreme of the disequilibrium structures the thesis studies.

\hypertarget{case-study-data}{%
\subsection{Case-study data}\label{case-study-data}}

The case study is built from the real Umpqua FORPLAN data (REF). The yield curves, i.e., the volume removed per acre by age for each ecoclass, prescription, and management intensity, come from the 1990 Final Environmental Impact Statement (FEIS) Appendix~B \citep{usda1990umpqua}. The managed yield tables in FORPLAN were generated with the DFSIM Douglas-fir simulator and adjusted to operational-falldown, while the mountain-hemlock yield curve was derived from Johnson's mountain-hemlock site equations. The economic parameters come from the FEIS's timber-value derivations: stumpage, logging, and manufacturing costs by species (Table~B-65), the Douglas-fir price-diameter and pond-value relationships (Table~B-66), and the per-ecoclass access (road) costs. These costs make the high-elevation mountain-hemlock (CM-CE) prescriptions negatively valued. The yield-curve shapes are consistent with the LRMP's documented 95\% CMAI culmination ages by ecoclass (Table~IV-3: 85~years for CH-CW, 120 for CD-CP, 115 for CR-CF, 175 for CM-CE), and the FEIS Stage~II financial analysis aligns with key values in the thesis: mature-timber NPV range of \$1,800--\$7,700 per acre, managed-prescription range of \$27--\$290 per acre, and the negative mountain-hemlock values.

All source documents were obtained and machine-read from HathiTrust (public domain): the 1990 LRMP, the 1990 FEIS main volume and its Appendix B (the FORPLAN analysis volume), and the 1987 Draft EIS \citep{usda1987umpqua} that Daugherty used. We performed a set of consistency checks to ensure a consistent reproduction of Daugherty's case-study data. First, the 1987 DEIS and 1990 FEIS yield tables agree, confirming that the same yield data were used. Second, the reconstructed case study reproduces the key values reported in Daugherty's Tables~5.3 and 5.4. The model NPVs of the five mature vegetation types match the Table~5.4 values exactly, while the managed prescriptions reproduce the signs and broad rotation ranges reported in Table 5.3. The exact ordering of optimal rotations is only partially reproduced because the FEIS prescription mapping simplifies some treatments (Section~5.5). Finally, we cross-checked the reconstructed mature volumes against the FEIS standing-volume curves at the corresponding ages and found them to be of the same order of magnitude. All reconstructed data and their provenance are documented in the repository (\texttt{src/fresh\\\_daugherty/instance/}; data summary and anchor tables in \texttt{supplementary/01-case-study-data.md}).

\hypertarget{the-sequential-replanning-simulator}{%
\subsection{The sequential-replanning simulator}\label{the-sequential-replanning-simulator}}

Dynamic inconsistency cannot be assessed from a single open-loop solution alone. It is revealed when the plan is subjected to sequential re-optimization. Following the thesis's iterative LP simulation, we implement a sequential-replanning simulator (Figure~\ref{fig:sequential-replanning}).

\begin{figure}[htbp]
\centering
\begin{tikzpicture}[
    node distance=0.5cm,
    >=latex,
    every node/.style={font=\small},
    box/.style={
        rectangle,
        rounded corners,
        draw,
        align=center,
        text width=8.5cm,
        minimum height=0.9cm
    },
    arrow/.style={->, thick}
]
\node[box] (solve) {
Solve the open-loop LP from the initial forest state over the planning horizon
};
\node[box, below=of solve] (apply) {
Apply the current-period decision from that solution
};
\node[box, below=of apply] (advance) {
Advance the forest state by one period: extract the realized area-by-age distribution and use it as the initial condition of the next subproblem
};
\node[box, below=of advance] (resolve) {
Re-solve the open-loop LP from the realized state over the rolling horizon
};
\node[box, below=of resolve] (record) {
Record the projected plan and the decision actually taken
};
\draw[arrow] (solve) -- (apply);
\draw[arrow] (apply) -- (advance);
\draw[arrow] (advance) -- (resolve);
\draw[arrow] (resolve) -- (record);
\draw[arrow]
    (record.east) -- ++(1.2,0)
    |- node[pos=0.25,right,align=left]{Repeat until the end\\of the simulation}
    (apply.east);
\end{tikzpicture}
\caption{Flow chart for the sequential-replanning simulation used to test dynamic consistency.}
\label{fig:sequential-replanning}
\end{figure}

Two modeling decisions bear on the interpretation. First, each replan re-solves the harvest-scheduling problem with a \emph{fresh harvest-flow constraint} anchored to that subproblem's own trajectory, not carried from the prior period's realized harvest. This \textit{rolling-horizon} convention is in line with FORPLAN practice, each decadal re-solve is a new plan with its own flow constraint, and it is what allows the originally promised non-declining flow to be abandoned. The alternative \textit{shrinking-horizon} convention (anchoring each replan's first period to the realized previous harvest, i.e., carrying the flow history) suppresses potential inconsistencies due to re-planning by forcing the realized trajectory to track the plan until the constraint becomes infeasible. We implemented it as an option and report the rolling- versus shrinking-horizon and reset conventions as a robustness consideration. Second, because LP optima can be non-unique, a divergence between the projected and realized harvest vectors could, in principle, reflect an alternate optimum rather than genuine inconsistency. To rule this out, we add an \emph{objective-gap diagnostic}: at each replan, we solve the subproblem freely and again with the period-1 harvest held to the announced plan's value (within a 1\% band), and we compare the two objectives. If the re-solved objective strictly exceeds the tail-fixed objective -- or the announced value is infeasible from the realized state -- the announced plan's tail is strictly suboptimal or unimplementable, i.e., dynamically inconsistent, not merely a different member of an optimal set.

The inconsistency measure is the divergence between the open-loop plan's projected per-period harvest and the realized replanned trajectory, i.e., the changes in decisions and in projected output levels over time between what was announced and what the re-solving planner actually does. Formally, with the open-loop projection \(p = (p_1, \ldots, p_T)\) and the realized replanned trajectory \(r = (r_1, \ldots, r_T)\), we measure the per-period divergence by the symmetric relative deviation
\begin{equation}
\delta_t = \frac{|p_t - r_t|}{\max(|p_t|, |r_t|, \varepsilon)},
\label{eq:metric}
\end{equation}
which is bounded in \([0, 1]\) and robust to near-zero harvest baselines, and summarizes a cell by its mean \(\bar\delta = \frac{1}{T}\sum_t \delta_t\) and by the relative change in total volume \((\sum_t r_t - \sum_t p_t)/\sum_t p_t\). A cell exhibits dynamic inconsistency when \(\bar\delta\) exceeds a small tolerance (5\%, well above LP-solver numerical noise and well below the observed magnitudes). The first-period decision is consistent by construction (the realized period-1 harvest is the open-loop period-1 decision), so the divergence is in the plan's tail, as the theory predicts.

\hypertarget{experiment-design}{%
\subsection{Scenario design}\label{experiment-design}}

The scenario grid reproduces the thesis's three sensitivity dimensions: initial forest condition (landbase), harvest policy, and the interest rate. The reproduced scenario grid spans \emph{all eighteen} of the thesis's initial forest conditions \citep[Table~5.5]{daugherty1991credibility}: the all-mature landbases (1, and 2 excluding the negatively valued CM-CE ecoclass); the area-control-derived landbases (3--8: landbase 1 or 2 after 40 or 70 years of area-control harvest at a 70- or 100-year rotation, a disequilibrium\(\to\)regulation gradient); the structured young-growth landbases (9, equal acres by age class; 10, unequal); and the randomly generated young-growth landbases (11--18). Harvest policy is the thesis's six harvest-flow constraint sets \citep[Table~5.6]{daugherty1991credibility}: no harvest flow constraint (NHF, the control); non-declining yield (NDY); bounded decline of 10\% or 20\% (\(-10\%\), \(-20\%\)); and symmetric bounded deviation of 10\% or 20\% (\(\pm10\%\), \(\pm20\%\)). The discount rate takes the thesis's four values (0\%, 2\%, 4\%, 6\%). Each of the resulting 432~scenarios is a full sequential-replanning simulation under the scenarios's policy (equation~\eqref{eq:openloop}), scored for the occurrence and magnitude of dynamic inconsistency (equation~\eqref{eq:metric}). The thesis re-solved the model for eleven periods over its planning horizon; we extend the simulation to the full fifteen-period horizon, increasing the observation window without changing the underlying mechanism.

The thesis also used ending-period (terminal) constraints with the ending inventory at least 80\% of the regulated forest's average inventory, and final-period harvest at most 120\% of long-term sustained yield to minimize horizon effects. We implemented and tested these and found that with the corrected case-study data, the model does not exhibit the horizon-end liquidation they guard against. These non-binding constraints were are omitted from the reported scenarios but remain available in the stack.

\hypertarget{extensions-four-mitigation-tests}{%
\subsection{Extensions: four mitigation experiments}\label{extensions-four-mitigation-tests}}

Beyond reproduction, we run four extension experiments (E1--E4), each replacing one ingredient of the model and asking the same question: does the change \emph{mitigate} or \emph{eliminate} dynamic inconsistency? Each extension is scored with the divergence metric (equation~\eqref{eq:metric}), occurrence criterion, and objective-gap diagnostic against the unchanged scenario grid as control. Detailed results, including the harvest schedules, are documented in the supplement (\texttt{supplementary/05}--\texttt{08}).

\textbf{E1 -- time-varying discount-rate shapes.} We replace the thesis's constant discount rate with time-varying, declining discount rates, considering two functional forms: linear decline (4\%$\to$0\% or 6\%$\to$0\% over the horizon) or an inverse-j shape (constant 4\% for the first one or two periods, then exponential decay halving the rate per period), with the objective's discount factors compounded period by period.  Declining discount-rate schedules weight the future more heavily, providing a candidate remedy if discounting-induced  front-loading contributes to dynamic inconsistency.

\textbf{E2 -- max-harvest-cap even-flow search.} The even-flow constraint is removed entirely and replaced by a per-period cap $H_t \leq \mathrm{cap}$, calibrated by bisection until until the realized sequential-replanning trajectory meets three even-flow criteria: trend slope within 1\% of the mean per period, maximum period-over-period fluctuation within 5\% of the mean, and coefficient of variation within 5\%). This is an automated allowable-cut calibration that searches for the highest harvest level that is both even and actually deliverable under replanning.

\textbf{E3 -- value-denominated flow constraints.} The bounded-deviation link is denominated in \emph{undiscounted net revenue} instead of volume (the flow rows carry $\sum \text{net price} \times \text{volume}$), keeping the six policy forms. Because the negatively valued CM-CE stratum \emph{subtracts} revenue, it can no longer serve as constraint-filler under a revenue floor. This implements a direct test of whether the negatively valued ``inconsistent variables'' are the cause of the phenomenon.

\textbf{E4 -- rolling-mean NDY.} The pointwise link is smoothed: each period's harvest is floored at the mean of the previous $k \in \{2, 3\}$ periods, $H_{t+1} \geq \frac{1}{k}\sum_{i=t-k+1}^{t} H_i$, under two anchoring institutions: \emph{within-plan} windows (each replan re-anchors freely, as in the core design) and \emph{realized-history} windows (the window reaches back into the realized past harvests, binding the future planner to what actually happened; periods where the floor cannot be sustained are recorded as relaxations).

\hypertarget{implementation-and-reproducibility}{%
\subsection{Implementation and reproducibility}\label{implementation-and-reproducibility}}

The complete stack -- typed, provenance-recorded case-study data; the \texttt{ws3} forest model; the explicit LP builder; the sequential-replanning simulator; the experiment runner; and the validation records -- is openly available (see Software availability). The thesis's bounded-deviation sequential-flow harvest-flow constraint (equation~\eqref{eq:openloop}) required a new even-flow constraint form in \texttt{ws3} (consecutive-period anchoring with separate decrease/increase tolerances), contributed upstream and released in \texttt{ws3} v1.1.0a5; it is a reusable addition for the classic FORPLAN harvest-flow policies. The package is test-backed, the environment is locked, and the simulation records needed to regenerate every number reported below are produced by the repository's experiment entry points.
